\documentclass{ctexart}
\pagestyle{empty} % 禁用页眉页脚，包括页码
\setlength{\parindent}{2em}
\usepackage[
  paperwidth=17cm,
  paperheight=25cm,
  margin=6pt
]{geometry}
\usepackage{titlesec}

% 修改 \section 字号为 14pt，加粗
\titleformat{\section}
  {\normalfont\bfseries\large} % 样式：\Large 可改为 \huge, \small 等
  {\thesection}{0em}{}         % 编号和间距设置
  \titlespacing{\section}{0pt}{10pt}{0pt}

\usepackage{multicol}  
\usepackage{amsmath}   


\begin{document}
\section*{Dormand-Prince RK45 方法}

Dormand-Prince法是一种可以自适应调整积分步长的显式Runge-Kutta方法，其同时估计4阶解与5阶解，其Butcher表如下：
\vspace{-3pt}
\[
    \begin{array}{c|ccccccc}
        c_i       & a_{i1}     & a_{i2}      & a_{i3}     & a_{i4}   & a_{i5}        & a_{i6}   & a_{i7} \\
        \hline
        0         &            &             &            &          &               &          &        \\
        1/5       & 1/5        &             &            &          &               &          &        \\
        3/10      & 3/40       & 9/40        &            &          &               &          &        \\
        4/5       & 44/45      & -56/15      & 32/9       &          &               &          &        \\
        8/9       & 19372/6561 & -25360/2187 & 64448/6561 & -212/729 &               &          &        \\
        1         & 9017/3168  & -355/33     & 46732/5247 & 49/176   & -5103/18656   &          &        \\
        1         & 35/384     & 0           & 500/1113   & 125/192  & -2187/6784    & 11/84    & 0      \\
        \hline
        b_i       & 35/384     & 0           & 500/1113   & 125/192  & -2187/6784    & 11/84    & 0      \\
        \hat{b}_i & 5179/57600 & 0           & 7571/16695 & 393/640  & -92097/339200 & 187/2100 & 1/40
    \end{array}
\]

RK45 的更新公式可统一写为：
\[
    \begin{aligned}
        k_i                        & = f\Big(t_n + c_i h,\, \boldsymbol{s}_n + h \sum_{j=1}^{i-1} a_{ij} k_j \Big), \quad i=1,\dots,7, \\
        \boldsymbol{s}_{n+1}^{(5)} & = \boldsymbol{s}_n + h \sum_{i=1}^{7} \hat{b}_i k_i,\quad
        \boldsymbol{s}_{n+1}^{(4)} = \boldsymbol{s}_n + h \sum_{i=1}^{7} b_i k_i                                                       \\
    \end{aligned}
\]


设$\varepsilon$为相对误差限，可使用下式对积分步长进行自适应调整：
\[
    h_{\text{new}} = h\cdot\left(\frac{\varepsilon h}{2\left|\boldsymbol{s}_{n+1}^{(5)}-\boldsymbol{s}_{n+1}^{(4)}\right|}\right)^{1/5}
\]
\rule{\linewidth}{0.4pt}
\begin{multicols}{2}
    \section*{二维弹跳小球状态空间模型}

    定义状态向量
    \begin{equation*}
        \boldsymbol{s} =
        \begin{bmatrix}
            x \\ y \\ v_x \\ v_y
        \end{bmatrix},
        \quad
        \dot{\boldsymbol{s}} =
        \begin{bmatrix}
            v_x \\ v_y \\ 0 \\ -g
        \end{bmatrix}
    \end{equation*}
    其中 $g$ 为重力加速度.

    \section*{碰撞模型（恢复系数）}

    记$e \in [0,1]$ 为恢复系数，碰撞时状态更新：
    \[
        v_y^+ = -e\, v_y^-
    \]

    \section*{碰撞模型（弹簧阻尼）}

    接触地面时竖直力：
    \[
        F_y = -k\,y - c\,v_y,
    \]

    状态空间表示：
    \[
        \dot{\boldsymbol{s}} =
        \begin{bmatrix}
            v_x \\ v_y \\ 0 \\ -g + F_y/m
        \end{bmatrix}, \quad y \le 0
    \]
\end{multicols}

\end{document}